% 🧪 Probe fixture of case scale-discrete, every scenario.
% Band and point scales are measured through the three answers d3 exposes — the start of each band,
% its width and the step between neighbours — because a family that draws a bar needs all three.
% The binned kinds are measured through the range entry they assign to a value.
\documentclass{article}
\usepackage{semio-viz-scale}
\usepackage{semio-viz-probe}
\ExplSyntaxOn
\fp_new:N \l_probe_lo_fp
\fp_new:N \l_probe_hi_fp
\fp_new:N \l_probe_value_fp
\seq_new:N \l_probe_seq
\tl_new:N \l_probe_item_tl

% 📦 Writes the left edge of every band of a band or point scale.
\cs_new_protected:Npn \probestarts #1#2#3 {
  \seq_clear:N \l_probe_seq
  \clist_map_inline:nn {#3}
    {
      \semio_viz_scale_band_edges:nnNN {#2} {##1} \l_probe_lo_fp \l_probe_hi_fp
      \seq_put_right:Nx \l_probe_seq { \fp_to_decimal:n { round( \l_probe_lo_fp , 6 ) } }
    }
  \semio_viz_probe_values:nx {#1} { \seq_use:Nn \l_probe_seq { , } }
}

% 📏 Writes the bandwidth and the step of a band or point scale.
\cs_new_protected:Npn \probemetrics #1#2 {
  \semio_viz_scale_band_width:nN {#2} \l_probe_lo_fp
  \semio_viz_scale_band_step:nN {#2} \l_probe_hi_fp
  \semio_viz_probe_values:nx {#1}
    {
      \fp_to_decimal:n { round( \l_probe_lo_fp , 6 ) } ,
      \fp_to_decimal:n { round( \l_probe_hi_fp , 6 ) }
    }
}

% 🎯 Writes the range entry a binned or ordinal scale assigns to each input.
\cs_new_protected:Npn \probeitems #1#2#3 {
  \seq_clear:N \l_probe_seq
  \clist_map_inline:nn {#3}
    {
      \semio_viz_scale_binned_item:nnN {#2} {##1} \l_probe_item_tl
      \seq_put_right:NV \l_probe_seq \l_probe_item_tl
    }
  \semio_viz_probe_values:nx {#1} { \seq_use:Nn \l_probe_seq { , } }
}

% 🔤 Writes the range entry an ordinal scale assigns to each category.
\cs_new_protected:Npn \probeordinal #1#2#3 {
  \seq_clear:N \l_probe_seq
  \clist_map_inline:nn {#3}
    {
      \semio_viz_scale_ordinal_item:nnN {#2} {##1} \l_probe_item_tl
      \seq_put_right:NV \l_probe_seq \l_probe_item_tl
    }
  \semio_viz_probe_values:nx {#1} { \seq_use:Nn \l_probe_seq { , } }
}

% 🪜 Writes the thresholds a binned scale derives from its domain and range.
\cs_new_protected:Npn \probethresholds #1#2 {
  \semio_viz_scale_thresholds:n {#2}
  \semio_viz_probe_values:nx {#1}
    { \seq_use:Nn \l_semio_viz_scale_thresholds_seq { , } }
}
\ExplSyntaxOff
\begin{document}

\SemioVizProbeBegin{scale-discrete}{band}
\SemioVizScale{plain}{band}{A,B,C,D,E}{0,100}
\SemioVizScale{padded}{band}{A,B,C,D}{0,120}[paddingInner=0.1,paddingOuter=0.2]
\SemioVizScale{aligned}{band}{A,B,C}{0,90}[paddingInner=0.5,paddingOuter=0.5,align=0]
\SemioVizScale{reversed}{band}{A,B,C,D}{120,0}[paddingInner=0.2]
\probestarts{band/plain}{plain}{A,B,C,D,E}
\probemetrics{metrics/plain}{plain}
\probestarts{band/padded}{padded}{A,B,C,D}
\probemetrics{metrics/padded}{padded}
\probestarts{band/aligned}{aligned}{A,B,C}
\probemetrics{metrics/aligned}{aligned}
\probestarts{band/reversed}{reversed}{A,B,C,D}
\probemetrics{metrics/reversed}{reversed}

\SemioVizProbeBegin{scale-discrete}{point}
\SemioVizScale{pt}{point}{A,B,C,D}{0,120}
\SemioVizScale{ptpad}{point}{A,B,C,D}{0,120}[padding=0.5]
\probestarts{point/pt}{pt}{A,B,C,D}
\probemetrics{metrics/pt}{pt}
\probestarts{point/ptpad}{ptpad}{A,B,C,D}
\probemetrics{metrics/ptpad}{ptpad}

\SemioVizProbeBegin{scale-discrete}{ordinal}
\SemioVizScale{ord}{ordinal}{A,B,C,D,E}{10,20,30}
\probeordinal{ordinal/ord}{ord}{A,B,C,D,E}

\SemioVizProbeBegin{scale-discrete}{quantize}
\SemioVizScale{qz}{quantize}{0,1}{1,2,3,4}
\SemioVizScale{qzwide}{quantize}{-10,10}{a,b,c}
\probeitems{quantize/qz}{qz}{-1,0,0.1,0.25,0.5,0.6,0.9,1,2}
\probethresholds{thresholds/qz}{qz}
\probeitems{quantize/qzwide}{qzwide}{-10,-4,0,4,10}
\probethresholds{thresholds/qzwide}{qzwide}

\SemioVizProbeBegin{scale-discrete}{quantile}
\SemioVizScale{ql}{quantile}{1,2,3,4,5,6,7,8,9,10}{10,20,30,40}
\probeitems{quantile/ql}{ql}{1,3,5,7,9,10}
\probethresholds{thresholds/ql}{ql}

\SemioVizProbeBegin{scale-discrete}{threshold}
\SemioVizScale{th}{threshold}{0,1}{low,mid,high}
\probeitems{threshold/th}{th}{-1,0,0.5,1,2}
\probethresholds{thresholds/th}{th}

\SemioVizProbePage
\end{document}
